\documentclass[11pt]{article}
\usepackage[margin=1in]{geometry}
\usepackage{amsmath,amssymb}
\setlength{\parindent}{0pt}
\setlength{\parskip}{0.7em}
\begin{document}
\textbf{21-127 Concepts of Mathematics (fictional)}

Homework 2 - Creative mathematical worlds --- Practice submission B

Student: Fictional learner

\section*{Question 1: The two-button robot}

Exactly all even integers can occur. Each button changes the display by an even integer, so the display is always even. Pressing P and then Q adds $2$. Repeating that pair reaches every even integer, which finishes the proof.

\newpage
\textbf{21-127 Concepts of Mathematics (fictional)}

Homework 2 - Creative mathematical worlds --- Practice submission B

Student: Fictional learner

\section*{Question 2: One key for every visitor?}

(a) is $\forall v\in V\ \exists k\in K\ R(v,k)$ and (b) is $\exists k\in K\ \forall v\in V\ R(v,k)$.

For a counterexample, let red work only for Ada, green only for Bo, and blue only for Cy. No other openings are allowed. Each visitor has a key, but every key fails for the other two visitors. Thus the implication is false.

\newpage
\textbf{21-127 Concepts of Mathematics (fictional)}

Homework 2 - Creative mathematical worlds --- Practice submission B

Student: Fictional learner

\section*{Question 3: Three self-referential cards}

A is true, B is false, and C is true. A correctly says that B is false, and B incorrectly says that C is false. Since those two cards agree with the assignment, this is a consistent solution. The cards force one another, so it is unique.

\newpage
\textbf{21-127 Concepts of Mathematics (fictional)}

Homework 2 - Creative mathematical worlds --- Practice submission B

Student: Fictional learner

\section*{Question 4: The archive's invert-selection button}

For each $x\in U$, $x\in T(T(A))$ means $x\notin U\setminus A$, hence $x\in A$. The equivalence goes both ways, and both sets are subsets of $U$. Thus $T(T(A))=A$.

I also considered whether half the files being selected could make a fixed point. Equal size does not imply equal sets. With any file $x\in U$, inversion changes whether that very file belongs, so when $U$ is nonempty the two selections cannot be equal.

For the empty archive there is just one selection, $A=\varnothing$, and it is fixed. This separate case is not a counterexample to the nonempty-archive claim.

\newpage
\textbf{21-127 Concepts of Mathematics (fictional)}

Homework 2 - Creative mathematical worlds --- Practice submission B

Student: Fictional learner

\section*{Question 5: The paired shelf labels}

Put the 13 chosen labels in increasing order. There are 12 gaps between consecutive chosen labels and 24 possible shelf positions. By the pigeonhole principle, one consecutive gap must be 12, so two chosen labels differ by exactly 12.

\end{document}
